\documentclass[12pt,a4paper]{article}
\usepackage[utf8]{inputenc}
\usepackage[T1]{fontenc}
\usepackage[english]{babel}
\usepackage{geometry}
\geometry{a4paper, margin=1in, top=0.8in, bottom=0.8in}

% Color and Graphics Packages
\usepackage[table]{xcolor}
\usepackage{graphicx}
\usepackage{tikz}
\usepackage{tcolorbox}
\usetikzlibrary{shapes,arrows,positioning,fit,backgrounds}

% Mathematics Packages
\usepackage{amsmath, amssymb, amsthm}
\usepackage{mathtools}
\usepackage{physics}
\usepackage{bm}

% Code and Listings
\usepackage{listings}
\usepackage{inconsolata}

% Table and Layout
\usepackage{array}
\usepackage{booktabs}
\usepackage{multirow}
\usepackage{multicol}
\usepackage{enumitem}

% Header Footer
\usepackage{fancyhdr}
\usepackage{hyperref}
\hypersetup{
    colorlinks=true,
    linkcolor=blue,
    urlcolor=blue,
}

% Custom Boxes
\usepackage{tcolorbox}
\tcbuselibrary{skins,breakable,listings}

% ==================== COLOR DEFINITIONS ====================
\definecolor{primary}{RGB}{25, 55, 95}
\definecolor{secondary}{RGB}{70, 130, 200}
\definecolor{accent}{RGB}{220, 70, 50}
\definecolor{success}{RGB}{50, 150, 80}
\definecolor{warning}{RGB}{200, 150, 30}
\definecolor{info}{RGB}{30, 150, 180}
\definecolor{codebg}{RGB}{245, 245, 245}
\definecolor{comment}{RGB}{100, 100, 100}
\definecolor{keyword}{RGB}{0, 0, 255}
\definecolor{string}{RGB}{163, 21, 21}
\definecolor{function}{RGB}{0, 134, 179}

% ==================== STYLE DEFINITIONS ====================
\lstset{
    language=Mathematica,
    backgroundcolor=\color{codebg},
    basicstyle=\ttfamily\small,
    keywordstyle=\color{keyword}\bfseries,
    commentstyle=\color{comment}\itshape,
    stringstyle=\color{string},
    showstringspaces=false,
    tabsize=2,
    breaklines=true,
    frame=single,
    framerule=0.5pt,
    rulecolor=\color{gray!30},
    numbers=left,
    numberstyle=\tiny\color{gray},
    numbersep=5pt,
}

% Custom TColorBoxes
\newtcolorbox{mainbox}[2][]{
    colback=white,
    colframe=primary,
    coltitle=white,
    fonttitle=\bfseries\large,
    title=#2,
    sharp corners,
    boxrule=1.5pt,
    left=8pt,
    right=8pt,
    top=8pt,
    bottom=8pt,
    #1
}

\newtcolorbox{commandbox}[1][]{
    colback=blue!5,
    colframe=secondary,
    sharp corners,
    boxrule=1pt,
    left=6pt,
    right=6pt,
    top=5pt,
    bottom=5pt,
    #1
}

\newtcolorbox{examplebox}[1][]{
    colback=green!5,
    colframe=success,
    sharp corners,
    boxrule=1pt,
    left=6pt,
    right=6pt,
    top=5pt,
    bottom=5pt,
    #1
}

\newtcolorbox{warningbox}[1][]{
    colback=yellow!5,
    colframe=warning,
    sharp corners,
    boxrule=1pt,
    left=6pt,
    right=6pt,
    top=5pt,
    bottom=5pt,
    #1
}

\newtcolorbox{tipbox}[1][]{
    colback=cyan!5,
    colframe=info,
    sharp corners,
    boxrule=1pt,
    left=6pt,
    right=6pt,
    top=5pt,
    bottom=5pt,
    #1
}

% Section Formatting
\usepackage{titlesec}
\titleformat{\section}
    {\color{primary}\normalfont\LARGE\bfseries}
    {\color{primary}\thesection}{1em}{}
\titleformat{\subsection}
    {\color{secondary}\normalfont\Large\bfseries}
    {\color{secondary}\thesubsection}{1em}{}
\titleformat{\subsubsection}
    {\color{accent}\normalfont\large\bfseries}
    {}{0em}{}

% Page Style
\pagestyle{fancy}
\fancyhf{}
\fancyhead[L]{\textcolor{primary}{\textbf{Mathematica Command Reference}}}
\fancyhead[R]{\textcolor{primary}{\textbf{Complete Guide}}}
\fancyfoot[C]{\thepage}
\fancyfoot[R]{\textcolor{gray}{Version 1.0}}

% Document
\begin{document}

% ==================== TITLE PAGE ====================
\begin{titlepage}
    \centering
    \vspace*{1cm}
    
    {\Huge\bfseries\color{primary} Mathematica}\\[0.3cm]
    {\Huge\bfseries\color{primary} Complete Command Reference}\\[0.5cm]
    
    \rule{0.7\textwidth}{1pt}\\[0.5cm]
    
    {\Large A Comprehensive Guide to Syntax, Examples, and Applications}\\[0.8cm]
    
    \begin{tikzpicture}[remember picture,overlay]
        \node[opacity=0.05] at (current page.center) {
            \begin{tikzpicture}
                \draw[->,thick] (0,0) -- (2,0);
                \draw[->,thick] (0,0) -- (0,2);
                \node at (2.2,0) {x};
                \node at (0,2.2) {y};
                \draw[domain=0:1.5,smooth,variable=\x,thick] 
                    plot ({\x}, {sin(deg(\x*2))*0.5});
            \end{tikzpicture}
        };
    \end{tikzpicture}
    
    \vspace{1cm}
    
    \begin{mainbox}[width=0.8\textwidth, center]{Quick Navigation}
        \begin{center}
            \begin{tabular}{c c c}
                \textbf{Lists \& Tables} & \textbf{Matrix Operations} & \textbf{Loops \& Control}\\
                \textbf{Calculus} & \textbf{Linear Algebra} & \textbf{Graphics}\\
                \textbf{Data Analysis} & \textbf{Strings} & \textbf{File I/O}
            \end{tabular}
        \end{center}
    \end{mainbox}
    
    \vfill
    
    {\large Department of Applied Mathematics}\\
    {\large University of Dhaka}\\
    \vspace{0.5cm}
    {\large \today}
    
\end{titlepage}

\tableofcontents
\newpage

% ==================== SECTION 1: LISTS AND TABLES ====================
\section{Lists and Tables}

\subsection{Basic List Creation}

\begin{commandbox}[Basic List Commands]
\begin{lstlisting}
{1, 2, 3, 4, 5}                    (* Simple list *)
Range[10]                           (* {1,2,3,4,5,6,7,8,9,10} *)
Range[0, 10, 2]                     (* {0,2,4,6,8,10} *)
ConstantArray[0, 5]                 (* {0,0,0,0,0} *)
\end{lstlisting}
\end{commandbox}

\subsection{Table Command - The Most Important}

\begin{mainbox}{Table Command}
\textbf{Syntax:} \texttt{Table[expression, \{iterator\}]}

\textbf{When to use:} Creating lists, matrices, or any repeated structure. More efficient than loops.

\begin{lstlisting}
(* 1D Table *)
Table[i^2, {i, 1, 5}]               (* {1,4,9,16,25} *)
Table[2^i, {i, 0, 5}]               (* {1,2,4,8,16,32} *)
Table[RandomInteger[10], {i, 1, 5}] (* Random list *)

(* 2D Table (Matrix) *)
Table[i*j, {i, 1, 3}, {j, 1, 3}]    
(* {{1,2,3},{2,4,6},{3,6,9}} *)

(* Table with condition *)
Table[If[EvenQ[i], i, Nothing], {i, 1, 10}]
(* {2,4,6,8,10} *)

(* Table with step *)
Table[i, {i, 0, 1, 0.1}]            (* {0,0.1,0.2,...,1} *)
\end{lstlisting}
\end{mainbox}

\begin{examplebox}[Real Example: Sine Values]
\begin{lstlisting}
(* Create table of sine values *)
sineTable = Table[{x, Sin[x]}, {x, 0, 2Pi, Pi/6}];
TableForm[sineTable, TableHeadings -> {None, {"Angle", "Sin"}}]
\end{lstlisting}
\end{examplebox}

\subsection{Grid and Display Commands}

\begin{commandbox}[Display Commands]
\begin{lstlisting}
(* TableForm - displays as table *)
TableForm[{{1,2,3},{4,5,6}}]

(* MatrixForm - with brackets *)
MatrixForm[{{1,2},{3,4}}]

(* Grid - most flexible *)
Grid[{{a,b,c},{d,e,f}}, Frame -> All]
Grid[{{"Name","Age"},{"Alice",25},{"Bob",30}}, 
     Dividers -> All]

(* Column and Row *)
Column[{a,b,c}]                     (* Vertical *)
Row[{a,b,c}, " | "]                 (* Horizontal with separators *)
\end{lstlisting}
\end{commandbox}

% ==================== SECTION 2: MATRIX OPERATIONS ====================
\section{Matrix Operations}

\subsection{Matrix Creation}

\begin{mainbox}{Special Matrices}
\begin{lstlisting}
IdentityMatrix[4]                    (* 4x4 identity *)
DiagonalMatrix[{1,2,3,4}]           (* Diagonal matrix *)
ConstantArray[0, {3,4}]             (* 3x4 zero matrix *)
RandomInteger[{-10,10}, {3,3}]      (* 3x3 random integers *)
RandomReal[{0,1}, {4,4}]            (* 4x4 random reals *)
\end{lstlisting}
\end{mainbox}

\subsection{Matrix Manipulation}

\begin{commandbox}[Basic Operations]
\begin{lstlisting}
A = {{1,2,3},{4,5,6},{7,8,9}};
B = {{9,8,7},{6,5,4},{3,2,1}};

A + B                               (* Addition *)
A - B                               (* Subtraction *)
3 * A                               (* Scalar multiplication *)
A . B                               (* Matrix multiplication (use DOT!) *)
A * B                               (* Element-wise multiplication *)

Transpose[A]                        (* Transpose *)
Inverse[A]                          (* Inverse (if exists) *)
Det[A]                              (* Determinant *)
Tr[A]                               (* Trace *)
MatrixRank[A]                       (* Rank *)
\end{lstlisting}
\end{commandbox}

\begin{warningbox}[Important Note]
Matrix multiplication uses \texttt{.} (dot), not \texttt{*} (star). The star performs element-wise multiplication!
\end{warningbox}

\subsection{Element Access and Submatrices}

\begin{commandbox}[Indexing]
\begin{lstlisting}
A = Table[i*j, {i,4}, {j,4}];

A[[2,3]]                            (* Element at row2, col3 *)
A[[2]]                              (* Second row *)
A[[All, 3]]                         (* Third column *)
A[[2;;4, 2;;3]]                     (* Rows 2-4, cols 2-3 *)

(* Submatrix extraction *)
Take[A, {1,3}, {2,4}]               (* Rows 1-3, cols 2-4 *)
Drop[A, {2}, {3}]                   (* Remove row2, col3 *)
\end{lstlisting}
\end{commandbox}

% ==================== SECTION 3: LOOPS AND CONTROL ====================
\section{Loops and Control Flow}

\subsection{Do Loop}

\begin{mainbox}{Do Loop}
\textbf{Syntax:} \texttt{Do[expression, \{iterator\}]}

\textbf{When to use:} When you want to repeat an operation without collecting results.

\begin{lstlisting}
(* Basic Do *)
Do[Print[i], {i, 1, 5}]              (* Prints 1,2,3,4,5 *)

(* Nested Do *)
Do[
  Do[Print[i, ",", j], {j, 1, 3}],
  {i, 1, 3}
]

(* Accumulation *)
sum = 0;
Do[sum += i^2, {i, 1, 100}];
Print[sum]                           (* Sum of squares *)
\end{lstlisting}
\end{mainbox}

\subsection{For Loop}

\begin{mainbox}{For Loop}
\textbf{Syntax:} \texttt{For[start, test, increment, body]}

\textbf{When to use:} Traditional loop when you need full control over initialization and increment.

\begin{lstlisting}
(* Basic For *)
For[i = 1, i <= 5, i++,
  Print[i]
]

(* For with step *)
For[i = 1, i <= 10, i += 2,
  Print[i]
]

(* For with break *)
For[i = 1, i <= 100, i++,
  If[i^2 > 1000, Break[]];
  Print[i^2]
]
\end{lstlisting}
\end{mainbox}

\subsection{While Loop}

\begin{mainbox}{While Loop}
\textbf{Syntax:} \texttt{While[condition, body]}

\textbf{When to use:} When the number of iterations is unknown.

\begin{lstlisting}
(* Newton's method example *)
x = 1.0;
While[Abs[x^2 - 2] > 0.0001,
  x = (x + 2/x)/2;
  Print[x]
]

(* Count digits *)
n = 12345;
count = 0;
While[n > 0,
  count++;
  n = Floor[n/10]
];
Print["Number of digits: ", count]
\end{lstlisting}
\end{mainbox}

\subsection{Conditional Statements}

\begin{commandbox}[If, Which, Switch]
\begin{lstlisting}
(* If statement *)
If[x > 0,
  Print["Positive"],
  Print["Negative or zero"]
]

(* Which - multiple conditions *)
Which[
  x < 0, "Negative",
  x == 0, "Zero",
  x > 0, "Positive"
]

(* Switch - pattern matching *)
Switch[day,
  "Mon", 1,
  "Tue", 2,
  "Wed", 3,
  _, "Unknown"
]
\end{lstlisting}
\end{commandbox}

% ==================== SECTION 4: CALCULUS ====================
\section{Calculus Commands}

\subsection{Differentiation}

\begin{mainbox}{D Command}
\textbf{Syntax:} \texttt{D[f, x]} or \texttt{D[f, \{x, n\}]} for nth derivative

\begin{lstlisting}
(* Basic derivatives *)
D[x^3, x]                           (* 3x^2 *)
D[Sin[x], x]                        (* Cos[x] *)
D[Log[x], x]                        (* 1/x *)
D[Exp[a*x], x]                      (* a E^(a x) *)

(* Higher order *)
D[x^5, {x, 2}]                      (* 20x^3 *)

(* Partial derivatives *)
D[x^2*y^3, x]                       (* 2x y^3 *)
D[x^2*y^3, y]                       (* 3x^2 y^2 *)

(* Function form *)
f[x_] := x^2;
f'[x]                               (* 2x *)
f''[x]                              (* 2 *)
\end{lstlisting}
\end{mainbox}

\subsection{Integration}

\begin{mainbox}{Integrate Command}
\textbf{Syntax:} \texttt{Integrate[f, x]} (indefinite) or \texttt{Integrate[f, \{x,a,b\}]} (definite)

\begin{lstlisting}
(* Indefinite integrals *)
Integrate[x^2, x]                   (* x^3/3 *)
Integrate[Sin[x], x]                (* -Cos[x] *)
Integrate[1/x, x]                   (* Log[x] *)

(* Definite integrals *)
Integrate[x^2, {x, 0, 3}]           (* 9 *)
Integrate[Sin[x], {x, 0, Pi}]       (* 2 *)

(* Multiple integrals *)
Integrate[x*y, {x,0,1}, {y,0,1}]    (* 1/4 *)

(* Improper integrals *)
Integrate[1/x^2, {x, 1, Infinity}]  (* 1 *)
Integrate[Exp[-x], {x,0,Infinity}]  (* 1 *)

(* Numerical integration *)
NIntegrate[Sin[x]/x, {x, 0, 1}]
\end{lstlisting}
\end{mainbox}

\subsection{Limits and Series}

\begin{commandbox}[Limits and Series]
\begin{lstlisting}
(* Limits *)
Limit[Sin[x]/x, x -> 0]             (* 1 *)
Limit[(1+1/n)^n, n -> Infinity]     (* E *)

(* One-sided limits *)
Limit[1/x, x -> 0, Direction -> "FromAbove"]  (* Infinity *)

(* Series expansion *)
Series[Sin[x], {x, 0, 5}]           (* Taylor series *)
Normal[Series[Cos[x], {x,0,4}]]     (* Convert to polynomial *)
\end{lstlisting}
\end{commandbox}

\subsection{Sums and Products}

\begin{commandbox}[Sum and Product]
\begin{lstlisting}
(* Summation *)
Sum[i^2, {i,1,10}]                  (* 385 *)
Sum[1/2^i, {i,0,Infinity}]          (* 2 *)
Sum[1/k^2, {k,1,Infinity}]          (* Pi^2/6 *)

(* Double sum *)
Sum[i*j, {i,1,3}, {j,1,3}]          (* 36 *)

(* Product *)
Product[i, {i,1,10}]                (* 3628800 = 10! *)
\end{lstlisting}
\end{commandbox}

% ==================== SECTION 5: LINEAR ALGEBRA ====================
\section{Linear Algebra}

\subsection{Vector Operations}

\begin{commandbox}[Vector Commands]
\begin{lstlisting}
v = {1, 2, 3};
w = {4, 5, 6};

v . w                               (* Dot product: 32 *)
Cross[v, w]                         (* Cross product: {-3,6,-3} *)
Norm[v]                             (* Magnitude: Sqrt[14] *)
Normalize[v]                        (* Unit vector *)
\end{lstlisting}
\end{commandbox}

\subsection{Eigenvalues and Eigenvectors}

\begin{mainbox}{Eigensystem}
\begin{lstlisting}
A = {{4,0,1},{2,3,2},{1,0,4}};

Eigenvalues[A]                      (* List of eigenvalues *)
Eigenvectors[A]                     (* List of eigenvectors *)
Eigensystem[A]                      (* Both {values, vectors} *)

CharacteristicPolynomial[A, x]      (* Characteristic polynomial *)

(* Eigenspace basis for specific eigenvalue *)
λ = 3;
NullSpace[A - λ*IdentityMatrix[3]]  (* Basis for eigenspace *)
\end{lstlisting}
\end{mainbox}

\subsection{Solving Linear Systems}

\begin{mainbox}{LinearSolve}
\begin{lstlisting}
A = {{1,1,2},{2,4,-3},{3,6,-5}};
b = {9,1,0};

(* Solve Ax = b *)
x = LinearSolve[A, b]               (* {3, -1, 2} *)

(* Using inverse *)
x = Inverse[A] . b

(* Cramer's Rule *)
detA = Det[A];
x1 = Det[ReplacePart[A, 1 -> b]]/detA;

(* Row reduction *)
augmented = Join[A, Transpose[{b}], 2];
RowReduce[augmented]
\end{lstlisting}
\end{mainbox}

\subsection{Rank and Null Space}

\begin{commandbox}[Rank and Nullity]
\begin{lstlisting}
MatrixRank[A]                       (* Rank of matrix *)
NullSpace[A]                        (* Basis for null space *)

(* Dimension Theorem: rank + nullity = number of columns *)
rank = MatrixRank[A];
nullity = Length[NullSpace[A]];
n = Dimensions[A][[2]];
rank + nullity == n                  (* True *)
\end{lstlisting}
\end{commandbox}

% ==================== SECTION 6: GRAPHICS ====================
\section{Graphics and Visualization}

\subsection{2D Plotting}

\begin{mainbox}{Plot Command}
\begin{lstlisting}
(* Basic plot *)
Plot[Sin[x], {x, 0, 2Pi}]

(* Multiple functions *)
Plot[{Sin[x], Cos[x]}, {x, 0, 2Pi},
  PlotLegends -> {"Sin", "Cos"}]

(* With styling *)
Plot[x^2, {x, -2, 2},
  PlotStyle -> {Red, Thick},
  AxesLabel -> {"x", "y"},
  PlotLabel -> "Parabola",
  GridLines -> Automatic,
  PlotRange -> {-1, 5}]

(* List plot *)
data = Table[{x, Sin[x]}, {x, 0, 2Pi, 0.1}];
ListPlot[data, Joined -> True]
\end{lstlisting}
\end{mainbox}

\subsection{Parametric and Polar Plots}

\begin{commandbox}[Special Plots]
\begin{lstlisting}
(* Parametric plot *)
ParametricPlot[{Cos[t], Sin[t]}, {t, 0, 2Pi}]

(* Polar plot *)
PolarPlot[1 + Cos[θ], {θ, 0, 2Pi}]

(* Contour plot *)
ContourPlot[x^2 + y^2, {x,-2,2}, {y,-2,2}]
\end{lstlisting}
\end{commandbox}

\subsection{3D Plotting}

\begin{commandbox}[3D Graphics]
\begin{lstlisting}
(* 3D surface *)
Plot3D[Sin[x]*Cos[y], {x,-2,2}, {y,-2,2}]

(* Parametric 3D curve *)
ParametricPlot3D[{Cos[t], Sin[t], t}, {t,0,4Pi}]

(* 3D with styling *)
Plot3D[x^2 - y^2, {x,-2,2}, {y,-2,2},
  ColorFunction -> "TemperatureMap",
  Mesh -> None,
  AxesLabel -> {"x","y","z"}]
\end{lstlisting}
\end{commandbox}

% ==================== SECTION 7: DATA ANALYSIS ====================
\section{Data Analysis}

\subsection{Statistical Functions}

\begin{commandbox}[Basic Statistics]
\begin{lstlisting}
data = {1,2,3,4,5,6,7,8,9,10};

Mean[data]                          (* Average: 5.5 *)
Median[data]                        (* Middle: 5.5 *)
Variance[data]                      (* Sample variance *)
StandardDeviation[data]             (* Standard deviation *)
Total[data]                         (* Sum: 55 *)
Min[data]                           (* Minimum: 1 *)
Max[data]                           (* Maximum: 10 *)

(* Descriptive statistics *)
Quantile[data, 0.25]                (* First quartile *)
\end{lstlisting}
\end{commandbox}

\subsection{Random Data Generation}

\begin{commandbox}[Random Data]
\begin{lstlisting}
(* Random integers *)
RandomInteger[10, 10]               (* 10 integers 0-10 *)
RandomInteger[{-10,10}, {4,4}]      (* 4x4 random matrix *)

(* Random reals *)
RandomReal[1, 5]                    (* 5 reals between 0-1 *)
RandomReal[{-1,1}, {3,3}]           (* 3x3 random matrix *)

(* Random choice *)
RandomChoice[{a,b,c,d}, 10]         (* 10 random picks *)

(* Random permutation *)
RandomSample[Range[10]]             (* Random ordering *)
\end{lstlisting}
\end{commandbox}

% ==================== SECTION 8: STRINGS ====================
\section{String Operations}

\begin{commandbox}[String Commands]
\begin{lstlisting}
s = "Hello World";

StringLength[s]                     (* 11 *)
StringJoin["Hello", " ", "World"]   (* "Hello World" *)
StringSplit[s]                      (* {"Hello","World"} *)
StringReplace[s, "World" -> "Mathematica"]
ToString[123.45]                    (* "123.45" *)
ToExpression["1 + 2 * 3"]           (* 7 *)

(* String formatting *)
StringForm["x = ``, y = ``", x, y]
\end{lstlisting}
\end{commandbox}

% ==================== SECTION 9: FILE OPERATIONS ====================
\section{File and System Operations}

\begin{commandbox}[File I/O]
\begin{lstlisting}
(* Working directory *)
Directory[]
SetDirectory["C:\\Data"]

(* File operations *)
FileNames[]                         (* List files *)
Import["data.csv"]                  (* Import data *)
Export["output.txt", data]          (* Export data *)

(* User input *)
name = Input["Enter your name: "]
x = Input["Enter a number: "]

(* Printing *)
Print["The value is ", x]
\end{lstlisting}
\end{commandbox}

% ==================== SECTION 10: UTILITY COMMANDS ====================
\section{Utility Commands}

\begin{commandbox}[Help and Information]
\begin{lstlisting}
?Sin                                (* Help for Sin *)
??Sin                               (* Detailed help *)
Names["Plot*"]                      (* All Plot functions *)
\end{lstlisting}
\end{commandbox}

\begin{commandbox}[Timing and Memory]
\begin{lstlisting}
Timing[Sum[i^2, {i,1,10000}]]       (* CPU time *)
AbsoluteTiming[expr]                (* Wall clock time *)
MemoryInUse[]                       (* Memory used *)
Clear[x, y, z]                      (* Clear variables *)
ClearAll["Global`*"]                (* Clear all variables *)
\end{lstlisting}
\end{commandbox}

\begin{commandbox}[System Information]
\begin{lstlisting}
$Version                           (* Mathematica version *)
$OperatingSystem                   (* OS info *)
$ProcessorCount                    (* Number of cores *)
Date[]                             (* Current date *)
TimeUsed[]                         (* CPU time used *)
\end{lstlisting}
\end{commandbox}

% ==================== SECTION 11: QUICK REFERENCE CARD ====================
\section{Quick Reference Card}

\begin{center}
\begin{tabular}{|l|l|}
\hline
\multicolumn{2}{|c|}{\textbf{\color{primary} MOST USED COMMANDS}} \\
\hline
\textbf{Category} & \textbf{Commands} \\
\hline
Lists & \texttt{Range, Table, List, Join, Append} \\
Matrix & \texttt{MatrixForm, Transpose, Inverse, Det} \\
Loops & \texttt{Do, For, While, Table} \\
Math & \texttt{Sin, Cos, Exp, Log, Sqrt, Abs} \\
Calculus & \texttt{D, Integrate, Limit, Series, Sum} \\
Linear Algebra & \texttt{LinearSolve, Eigenvalues, NullSpace, Rank} \\
Graphics & \texttt{Plot, Plot3D, ParametricPlot, ListPlot} \\
Statistics & \texttt{Mean, Median, Variance, RandomInteger} \\
Conditionals & \texttt{If, Which, Switch} \\
Utility & \texttt{Print, Clear, Timing, ?} \\
\hline
\end{tabular}
\end{center}

% ==================== SECTION 12: COMMON PATTERNS ====================
\section{Common Programming Patterns}

\begin{commandbox}[Functional Programming]
\begin{lstlisting}
(* Map - apply function to each element *)
Map[f, list]                        (* {f[1], f[2], ...} *)
f /@ list                           (* Same as Map *)

(* Apply - use function on whole list *)
Apply[Plus, list]                   (* Sum all elements *)
Plus @@ list                        (* Same as Apply *)

(* Select - filter list *)
Select[list, EvenQ]                 (* Keep even numbers *)

(* Fold - accumulate left to right *)
Fold[f, init, list]                 (* f[f[f[init,x1],x2],...] *)

(* Nest - apply repeatedly *)
Nest[f, x, n]                       (* f[f[...[x]...]] n times *)
\end{lstlisting}
\end{commandbox}

\begin{tipbox}[Performance Tips]
\begin{itemize}
    \item Use \texttt{Table} instead of \texttt{Do} + \texttt{AppendTo} for list building
    \item Use \texttt{Total} instead of \texttt{Sum} for large lists
    \item Use \texttt{Map} instead of \texttt{For} loops for element-wise operations
    \item Use \texttt{LinearSolve} instead of \texttt{Inverse} for solving systems
    \item Use \texttt{NIntegrate} for numerical integration of complicated functions
\end{itemize}
\end{tipbox}

\begin{warningbox}[Common Mistakes]
\begin{itemize}
    \item Matrix multiplication: use \texttt{.} not \texttt{*}
    \item Indexing starts at 1, not 0
    \item \texttt{==} for equality, \texttt{=} for assignment
    \item \texttt{;} suppresses output, no semicolon prints result
    \item Functions use square brackets: \texttt{f[x]}, not \texttt{f(x)}
\end{itemize}
\end{warningbox}

% ==================== END OF DOCUMENT ====================
\end{document}
